\section{Plan de trabajo y actividades pendientes}\label{sec:pendientes}

\newcommand{\hecho}{\ding{51}}
\newcommand{\falta}{\ding{113}}

\subsection{Preparación}
\begin{itemize}[label={}]
  \item \hecho{} Definir la pregunta y el problema.
  \item \hecho{} Identificar las fuentes de información.
  \item \hecho{} Crear el entorno de Python (\archivo{.venv}) y el notebook de contexto (\archivo{notebooks/00_contexto_y_datasets.ipynb}).
  \item \falta{} Completar la estructura de carpetas (\archivo{src/}, \archivo{docs/}, \archivo{data/interim/}, \archivo{data/processed/}).
  \item \falta{} Crear el control de versiones y el archivo \archivo{data/raw/MANIFEST.csv}.
  \item \falta{} Confirmar los cortes de los periodos de análisis.
\end{itemize}

\subsection{Descarga (14 conjuntos seleccionados)}
\begin{itemize}[label={}]
  \item \hecho{} Seleccionar los conjuntos de datos (Tabla~\ref{tab:seleccion}).
  \item \hecho{} Scripts de descarga con MANIFEST (Sección~\ref{sec:descarga}).
  \item \hecho{} Movilidad: Metro, Metrobús, Ecobici, hechos de tránsito (comparable, ampliada y 2024).
  \item \hecho{} Seguridad: carpetas FGJ 2016--2024, llamadas al 911 2019--2022.
  \item \falta{} Seguridad: SESNSP municipal (descarga manual).
  \item \hecho{} Economía: DENUE, ocupación hotelera, ENOE (35 trimestres), IMSS (128 meses).
  \item \hecho{} Pandemia y referencia: casos COVID, semáforo federal, Censo 2020, Marco Geoestadístico.
\end{itemize}

\subsection{Tareas surgidas en la descarga}
\begin{itemize}[label={}]
  \item \falta{} Unificar los formatos de fecha del 911 y revisar los traslapes entre semestres.
  \item \falta{} Documentar el cambio de registro de los hechos de tránsito (2021--2022).
  \item \falta{} Fijar la ventana post-pandemia común (2023-01 -- 2024-07).
  \item \falta{} Rellenar las semanas 2020-W53 y 2022-W01 del semáforo.
\end{itemize}

\subsection{Exploración}
\begin{itemize}[label={}]
  \item \falta{} Obtener el perfil de cada conjunto: filas, columnas, rango de fechas, \% de nulos, tipos y categorías.
  \item \falta{} Documentar los cambios metodológicos de cada fuente.
\end{itemize}

\subsection{Limpieza (selección final)}
\begin{itemize}[label={}]
  \item \hecho{} Catálogo de alcaldías (clave INEGI y polígonos) y normalización de texto.
  \item \hecho{} Carpetas FGJ, Metrobús e IMSS.
  \item \hecho{} Metro: codificación, nombres, duplicados, tipos de ceros y cierres de las Líneas 1 y 12 marcados.
  \item \falta{} Metro: usar en la integración el promedio diario y una variante sin las Líneas 1 y 12.
  \item \falta{} Semáforo: rellenar 2020-W53 y 2022-W01.
  \item \hecho{} Hechos de tránsito (mensual y diario), ocupación hotelera y casos COVID (diario, semanal y mensual).
  \item \falta{} Ocupación hotelera: \emph{backcasting} 2016--2018.
  \item \falta{} Población por alcaldía (Censo 2020) para calcular tasas.
\end{itemize}

\subsection{Integración y entrega}
\begin{itemize}[label={}]
  \item \falta{} Construir las bases \archivo{cdmx\_mensual\_alcaldia} y \archivo{cdmx\_semanal}.
  \item \falta{} Calcular las tasas y los índices, y añadir las variables de periodo y semáforo.
  \item \falta{} Validar contra las cifras oficiales y elaborar el diccionario de datos.
  \item \falta{} Hacer un análisis descriptivo con gráficas antes, durante y después de la pandemia.
\end{itemize}
